\centerline{\bf \Huge }

\centerline{\bf \Huge 1 билет}
\centerline{\bf \Huge }

1. Ромб. Признаки ромба. Свойства ромба.

2. Вывод площади трапеции.

3. Элементы прямоугольного треугольника. Признаки равенства прямоугольных треугольников. Теорема о медиане прямоугольного треугольника. Прямоугольный треугольник с углом 30°.

4. В четырехугольнике $MPKH$ диагонали пересекаются в точке $O, \angle O M H=\angle O H M$,

$P H=M K, P K=M H$. Найдите угол $M H K$.


5. В трапеции $MPHK, MK$ - большее основание. Площади треугольников $M H K$ и $K H P$ равны $S_1$ и $S_2$ соответственно. Найдите площадь трапеции.


6. В треугольнике $A B C$ на сторонах $A B$ и $B C$ выбраны соответственно точки $D$ и $E$. $C D$ и $A E$ пересекаются в точке $O, S_{A D O}=S_{C E O}=8, S_{D O E}=4$. Найдите площадь треугольника $A B C$.






\centerline{\bf \Huge }
\centerline{\bf \Huge 2 билет}
\centerline{\bf \Huge }

1. Параллелограмм. Признаки параллелограмма. Свойства параллелограмма.

2. Доказательство теоремы Пифагора.

3. Средняя линия треугольника. Её свойства и признаки.

4. В ромбе $M P H K$ диагонали пересекаются в точке $O$. На сторонах $MK, KH, PH$ взяты точки $A, B, C$ соответственно, $A K=K B=P C$. Докажите, что $O A=O B$, и найдите сумму углов $POC$ и $MOA$.

5. В трапеции $A B C D$, $AD$ - большее основание, $\angle D=60^{\circ}$. Биссектрисы углов $C$ и $D$ пересекаются в точке $O$, $O D=a, B C=b, A D=c$. Найдите площадь трапеции.

6. Внутри треугольника $A B C$ взята точка $D$ и соединена с его вершинами $A$ и $B$. На стороне $B C$ вне его построен треугольник $B C E$ так, что $\angle E B C=\angle A B D$ и $\angle E C B$ $=\angle B A D$. Докажите, что $\triangle D B E \sim \triangle A B C$.





\centerline{\bf \Huge }
\centerline{\bf \Huge 3 билет}
\centerline{\bf \Huge }

1. Прямоугольник. Признаки прямоугольника. Свойства прямоугольника.

2. Доказательство теоремы Фалеса (обратная, обобщенная)

3. Окружность. Все объекты связанные с окружностью и их свойства.

4. В прямоугольнике $MPKH$, $O$ - точка пересечения диагоналей, $PA$ и $H B$ - перпендикуляры, проведенные из вершин $P$ и $H$ к прямой $M K$. Известно, что $M A=O B$. Найдите угол РОМ.

5. В треугольнике $A B C, \angle C=90^{\circ}$. На сторонах $A C, A B, B C$ соответственно взяты точки $M, P, K$ так, что четырехугольник $C M P K$ является квадратом, $A C=6 \mathrm{~cm}, B C=$ $14 \mathrm{~cm}$. Найдите сторону $MC$.


6. В параллелограмме $A B C D ,\angle A$ острый, $C E \perp A B, B C=$ $2 A B, M$ - середина $A D$. Докажите, что $\angle E M D=$ $3 \angle A E M$.





\centerline{\bf \Huge }
\centerline{\bf \Huge 4 билет}
\centerline{\bf \Huge }

1. Квадрат. Признаки квадрата. Свойства квадрата.

2. Вывод площади треугольника.

3. Окружность. Все объекты связанные с окружностью и их свойства.

4. В треугольнике $M P K, \angle M=90^{\circ}, M P=M K$. На сторонах $M P, P K, M K$ отмечены точки $A, B, C$ соответственно так, что четырехугольник $M A B C$ является квадратом, $A C=a$. Найдите $P K$.

5. Внутри параллелограмма $A B C D$ отмечена точка $M$. Докажите, что сумма площадей треугольников $A M D$ и $BMC$ равна половине площади параллелограмма.

6. В трапеции $A B C D, A C \perp B D, C E \perp A D, A C=15 \mathrm{~cm}$, $A E=9 \mathrm{cm}$. Найдите площадь трапеции.


\newpage
\centerline{\bf \Huge }
\centerline{\bf \Huge 5 билет}
\centerline{\bf \Huge }

1. Трапеция. Признаки трапеции. 

2. Вывод площади параллелограмма.

3. Точка пересечения медиан треугольника. Свойства подобных треугольников. Признаки равнобедренного треугольника. 

4. В трапеции $A B C D$, $AD$ - большее основание. Через середину стороны $C D$ и вершину $B$ проведена прямая, пересекающая луч $A D$ в точке $E$. Докажите, что площадь трапеции равна площади треугольника $A B E$.

5. Треугольник $A B C$ прямоугольный $\left(\angle C=90^{\circ}\right), D \in A C$ и $E \in A B$, причем $D E \| B C$ и $D E=D C, A E=15$ мм, $E B$ $=20$ мм. Найдите периметр треугольника $A B C$.


6. Отрезок $B K$ ( $K$ принадлежит стороне $A C$ ) разбивает треугольник $A B C$ на два подобных треугольника $A B K$ и $K B C$, причем $\dfrac{S_{A B K}}{S_{B K C}}=\dfrac{1}{3}$. Найдите углы треугольника.

\newpage
\centerline{\bf \Huge }
\centerline{\bf \Huge 6 билет}
\centerline{\bf \Huge }

1. Свойства и признаки равнобокой трапеции.

2. Определение синуса, косинуса, тангенса и котангенса острого угла
прямоугольного треугольника. Все их свойства. Синус, косинус, тангенс, котангенс 0°, 30°,  45°, 60°, 90°

3. Теорема Вариньона.

4. Биссектриса угла $B$ прямоугольника $A B C D$ пересекает сторону $A D$ в точке $K, A K=5 \mathrm{~cm}, K D=7 \mathrm{~cm}$. Найдите площадь прямоугольника.

5. Даны два отрезка $A B$ и $E K$. Точки $C$ и $M$ лежат соответственно на отрезках $A B$ и $E K$. Отрезки $A C, C B$ и $E M, M K$ пропорциональны. Докажите, что $A B \cdot M K=$ $=C B \cdot E K$.

6. В треугольнике $A B C, C M$ - биссектриса внешнего угла $C, B M \perp C M$. Угол $B$ в два раза меньше этого внешнего угла, $B F \perp A M, A F: F M=3: 1$. Найдите угол $B A M$.

